% experiments_revised.tex -- Experimental Evaluation, reframed around the question the data answers.
% Numbers are from the sealed per-seed rows (n=128 per arm). Deployed win rates and the
% worst-pole summary were recomputed from the same rows on 2026-10-01.

\section{Experimental Evaluation}
\label{sec:experiments}

\ADcomment{KB, add your draft description of the results with the key numbers
(mean + std) in latex tables. I will take it from there. This includes baseline
numbers as well as numbers for cases with different types of noises with 2v2,
4v4 and 6v6. Claude can easily create those latex tables with proper captions +
implications of the results.}

The evaluation asks three questions:
\begin{enumerate}
  \item Does the role-allocated system hold two complementary strategies at
  every team size (\(\Delta_A>0\) and \(\Delta_B>0\))?
  \item Do the natural alternatives---separate specialists, parameter-shared
  students, or a single generalist---achieve the same, and what does each one
  cost or gain in raw win rate?
  \item Does the specialization survive deployment noise?
\end{enumerate}
All evaluations use \(n=128\) seeds per arm. Tables report mean \(\pm\)
standard deviation of per-seed paired win differences, which take values in
\(\{-1,0,+1\}\); the standard deviation is therefore large relative to the
mean, and confidence intervals are shown in the figures.

\subsection{Specialization Across Team Sizes}

\begin{table}[t]
\centering
\caption{Specialization of the role-allocated system. Mean $\pm$ std across
128 confirmatory seeds. Positive $\Delta_A$ and $\Delta_B$ mean each strategy
beats the other on its own regime.}
\label{tab:main_specialization}
\footnotesize
\setlength{\tabcolsep}{4pt}
\begin{tabular}{c|cc}
\hline
Team size & $\Delta_A$ & $\Delta_B$ \\
\hline
2v2 & $+0.117 \pm 0.749$ & $+0.539 \pm 0.546$ \\
4v4 & $+0.156 \pm 0.657$ & $+0.406 \pm 0.620$ \\
6v6 & $+0.062 \pm 0.411$ & $+0.117 \pm 0.346$ \\
\hline
\end{tabular}
\end{table}

Both effects are positive in mean at all three team sizes
(Table~\ref{tab:main_specialization}). Separation is largest at 2v2 and 4v4.
At 6v6 the binary effects are smaller (\(+0.062\), \(+0.117\)), largely
because every evaluated system wins 0.86--0.98 of 6v6 games, which leaves
little room for win/loss outcomes to differ. On score margin, which is not
capped in this way, the 6v6 strategies separate clearly: the
intended-direction margin advantage is \(+0.812\) on Pole~A and \(+0.734\)
on Pole~B, with both bootstrap intervals above zero. An earlier sealed
\(n=64\) exploratory 6v6 evaluation showed the same pattern.

\subsection{Baseline Comparison}

\paragraph{Summary.}
Table~\ref{tab:scorecard} condenses the baseline results into the quantity
the method is designed to deliver: the specialization on the \emph{weaker}
pole, \(\min(\Delta_A,\Delta_B)\). A system is only useful for
opponent-conditioned play if both strategies are distinct, so the weaker
pole is the binding one. We introduce this summary for readability; it is
computed from the same rows as the per-pole tables below and was not a
preregistered criterion.

\begin{table}[t]
\centering
\caption{Worst-pole specialization $\min(\Delta_A,\Delta_B)$ (mean, 128 seeds
per arm). Negative values mean one strategy dominates the other on both
regimes. Ours is the only system that is positive at every team size, and the
highest at every team size.}
\label{tab:scorecard}
\footnotesize
\setlength{\tabcolsep}{5pt}
\begin{tabular}{@{}lccc@{}}
\hline
Method & 2v2 & 4v4 & 6v6 \\
\hline
Specialists (no roles) & $-0.055$ & $-0.156$ & $+0.031$ \\
Share-Encoder          & $+0.070$ & $-0.156$ & $+0.039$ \\
Fully Shared+$z$       & $+0.063$ & $-0.109$ & $+0.008$ \\
Ours (heuristic roles) & $\mathbf{+0.117}$ & $\mathbf{+0.156}$ & $\mathbf{+0.062}$ \\
\hline
\end{tabular}
\end{table}

The per-pole picture explains this. In every baseline the deficit is on
Pole~B: \(\Delta_B\) is negative or close to zero
(Table~\ref{tab:sharing_baselines}). The cause is visible in the raw win
rates. Without explicit roles, the Pole-A specialist is a strong all-round
policy: at 4v4 it wins \(0.812\) on Pole~A and \(0.742\) on Pole~B, more
than the Pole-B specialist wins on its own regime (\(0.586\)). There is then
no reason to ever select the Pole-B strategy, and the ``two strategies''
collapse into one. Parameter sharing does not fix this: the shared students
copy the structure of their source policies, including the dominance.

\paragraph{Role organization creates the missing distinction.}
Committing \(k\) defenders turns the Pole-A strategy into an explicit
attack-heavy posture that is strong on Pole~A and weak on Pole~B, which is
exactly what makes the choice of strategy matter
(Table~\ref{tab:role_baseline_comparison}).
At 2v2 (matched seeds), \(\Delta_B\) rises from \(-0.055\) to \(+0.523\), and
the Pole-A strategy's win rate on Pole~B falls from \(0.648\) to \(0.070\).
At 4v4, \(\Delta_B\) changes from \(-0.156\) to \(+0.406\), and the Pole-A
win rate on Pole~B falls from \(0.742\) to \(0.102\) (separate seed blocks;
descriptive). At 6v6 (matched seeds), \(\Delta_B\) rises from \(+0.031\) to
\(+0.117\); the paired change is \(+0.086\) with 95\% interval
\([+0.023,+0.148]\), and the Pole-A win rate on Pole~B falls from \(0.945\)
to \(0.859\).

\begin{table}[t]
\centering
\caption{Effect of heuristic role allocation. Mean $\pm$ std across 128
seeds. 2v2 and 6v6 rows within a team size use matched seed blocks; 4v4 rows
come from separate sealed blocks.}
\label{tab:role_baseline_comparison}
\footnotesize
\setlength{\tabcolsep}{3pt}
\begin{tabular}{@{}clcc@{}}
\hline
Team & Method & $\Delta_A$ & $\Delta_B$ \\
\hline
2v2 & Specialists (no roles) & $+0.055 \pm 0.619$ & $-0.055 \pm 0.668$ \\
    & Ours (heuristic roles) & $+0.180 \pm 0.681$ & $+0.523 \pm 0.546$ \\
\hline
4v4 & Specialists (no roles) & $+0.273 \pm 0.611$ & $-0.156 \pm 0.620$ \\
    & Ours (heuristic roles) & $+0.156 \pm 0.657$ & $+0.406 \pm 0.620$ \\
\hline
6v6 & Specialists (no roles) & $+0.094 \pm 0.386$ & $+0.031 \pm 0.215$ \\
    & Ours (heuristic roles) & $+0.062 \pm 0.411$ & $+0.117 \pm 0.346$ \\
\hline
\end{tabular}
\end{table}

\begin{table}[t]
\centering
\caption{Parameter-sharing baselines distilled from the same teachers and
dataset. Mean $\pm$ std across 128 seeds. Sharing reproduces the structure of
the source policies rather than creating new specialization.}
\label{tab:sharing_baselines}
\footnotesize
\setlength{\tabcolsep}{3pt}
\begin{tabular}{@{}clcc@{}}
\hline
Team & Method & $\Delta_A$ & $\Delta_B$ \\
\hline
2v2 & Specialists (no roles) & $+0.055 \pm 0.619$ & $-0.055 \pm 0.668$ \\
    & Share-Encoder          & $+0.125 \pm 0.652$ & $+0.070 \pm 0.678$ \\
    & Fully Shared+$z$       & $+0.078 \pm 0.635$ & $+0.063 \pm 0.661$ \\
    & Ours (heuristic roles) & $+0.117 \pm 0.749$ & $+0.539 \pm 0.546$ \\
\hline
4v4 & Specialists (no roles) & $+0.273 \pm 0.611$ & $-0.156 \pm 0.620$ \\
    & Share-Encoder          & $+0.281 \pm 0.627$ & $-0.156 \pm 0.633$ \\
    & Fully Shared+$z$       & $+0.367 \pm 0.545$ & $-0.109 \pm 0.655$ \\
    & Ours (heuristic roles) & $+0.156 \pm 0.657$ & $+0.406 \pm 0.620$ \\
\hline
6v6 & Specialists (no roles) & $+0.094 \pm 0.386$ & $+0.031 \pm 0.215$ \\
    & Share-Encoder          & $+0.039 \pm 0.341$ & $+0.039 \pm 0.292$ \\
    & Fully Shared+$z$       & $+0.031 \pm 0.396$ & $+0.008 \pm 0.235$ \\
    & Ours (heuristic roles) & $+0.062 \pm 0.411$ & $+0.117 \pm 0.346$ \\
\hline
\end{tabular}
\end{table}

\paragraph{What specialization costs in raw win rate.}
Specialization is not free, and its cost depends on team size.
Table~\ref{tab:deployed_wr} reports the win rate each system obtains when
deployed as intended: its Pole-A strategy against Pole~A and its Pole-B
strategy against Pole~B (the Generalist plays its single policy on both).
At 2v2, ours has the highest Pole-A win rate (\(0.641\)) and is within
\(0.055\) of the best Pole-B win rate (Share-Encoder, \(0.641\)). At 6v6, ours is within \(0.03\) of the best baseline
on both regimes and above the Generalist on both. At 4v4, however, ours has
the \emph{lowest} win rate on both regimes: committing two of four agents to
defense costs \(0.11\)--\(0.17\) on Pole~A and \(0.02\)--\(0.19\) on Pole~B
relative to the baselines. The 4v4 baselines win more often precisely
because they use one dominant attack-capable policy on both regimes; they
cannot offer a second behavior.

\begin{table}[t]
\centering
\caption{Deployed win rate: each system's Pole-A strategy on Pole~A
(WR$_A$) and Pole-B strategy on Pole~B (WR$_B$); the Generalist plays its
single policy. Means over 128 seeds. Arms at a team size share the
evaluation protocol; seed blocks differ between the role-allocated row and
the sharing rows.}
\label{tab:deployed_wr}
\footnotesize
\setlength{\tabcolsep}{3pt}
\begin{tabular}{@{}lcccccc@{}}
\hline
& \multicolumn{2}{c}{2v2} & \multicolumn{2}{c}{4v4} & \multicolumn{2}{c}{6v6} \\
Method & WR$_A$ & WR$_B$ & WR$_A$ & WR$_B$ & WR$_A$ & WR$_B$ \\
\hline
Generalist             & 0.508 & 0.594 & 0.758 & 0.695 & 0.859 & 0.945 \\
Specialists (no roles) & 0.523 & 0.594 & 0.812 & 0.586 & 0.953 & 0.977 \\
Share-Encoder          & 0.555 & 0.641 & 0.805 & 0.523 & 0.938 & 0.977 \\
Fully Shared+$z$       & 0.477 & 0.586 & 0.820 & 0.547 & 0.898 & 0.953 \\
Ours (heuristic roles) & 0.641 & 0.586 & 0.648 & 0.508 & 0.922 & 0.977 \\
\hline
\end{tabular}
\end{table}

\paragraph{Generalist.}
Because the Generalist has no strategy label, it has no \(\Delta_A\) or
\(\Delta_B\); we compare it through the matched-seed gap
\(\Delta_{G,j}=V(\mathrm{ours},j)-V(\pi_G,j)\)
(Table~\ref{tab:generalist_baseline}). Ours is ahead on Pole~A at 2v2
(\(+0.195\)), behind on both regimes at 4v4 (\(-0.070\) each), and slightly
ahead on both at 6v6 (\(+0.039\), \(+0.031\)). The Generalist shows that a
single policy can reach high win rates; it cannot be steered toward a
different behavior when the opponent changes.

\begin{table}[t]
\centering
\caption{Generalist baseline. Win rate of the single policy on each regime,
and the matched-seed gap to the role-allocated system (mean $\pm$ std across
128 seeds).}
\label{tab:generalist_baseline}
\footnotesize
\setlength{\tabcolsep}{3pt}
\begin{tabular}{@{}ccccc@{}}
\hline
Team & WR$_A(\pi_G)$ & WR$_B(\pi_G)$ & $\Delta_{G,A}$ & $\Delta_{G,B}$ \\
\hline
2v2 & 0.508 & 0.594 & $+0.195 \pm 0.722$ & $-0.023 \pm 0.693$ \\
4v4 & 0.758 & 0.695 & $-0.070 \pm 0.654$ & $-0.070 \pm 0.604$ \\
6v6 & 0.859 & 0.945 & $+0.039 \pm 0.442$ & $+0.031 \pm 0.175$ \\
\hline
\end{tabular}
\end{table}

\paragraph{When to prefer the role-allocated system.}
The data support a specific claim. If the requirement is a team with two
genuinely different, selectable strategies, the role-allocated system is the
only one of the four two-strategy systems that delivers it at every team
size, with the largest
\(\Delta_B\) and the largest worst-pole specialization at 2v2, 4v4, and 6v6.
Training specialists separately does not produce complementary strategies
(one policy dominates), and parameter sharing inherits that dominance. If
the requirement is maximum win rate against these two fixed regimes, a
dominant single policy is competitive and, at 4v4, better. The
role-allocated system is therefore the appropriate choice when the
strategy itself must be controllable---for example, when a higher-level
selector or operator chooses the team's posture---and its cost in raw win
rate is scale-dependent: a gain at 2v2, a loss at 4v4, and roughly even at
6v6.

\subsection{Robustness to Deployment Noise}

We perturb only evaluation: localization noise, motion error, and a
two-decision control delay, each at medium severity, on the final
role-allocated system. Each team size uses one matched seed block for the
nominal and the three perturbed conditions.

\begin{table*}[t]
\centering
\caption{Specialization under deployment-time perturbations (medium
severity), mean $\pm$ std across 128 matched seeds. Nominal runs on the same
blocks: 2v2 $+0.219/+0.492$, 4v4 $+0.234/+0.484$, 6v6 $+0.031/+0.094$
($\Delta_A/\Delta_B$).}
\label{tab:noise_results}
\footnotesize
\setlength{\tabcolsep}{5pt}
\renewcommand{\arraystretch}{1.08}
\begin{tabular}{c|cc|cc|cc}
\hline
& \multicolumn{2}{c|}{Localization}
& \multicolumn{2}{c|}{Motion}
& \multicolumn{2}{c}{Delay} \\
Team & $\Delta_A$ & $\Delta_B$ & $\Delta_A$ & $\Delta_B$ & $\Delta_A$ & $\Delta_B$ \\
\hline
2v2 & $+0.203 \pm 0.644$ & $+0.531 \pm 0.560$
    & $+0.234 \pm 0.693$ & $+0.539 \pm 0.546$
    & $+0.258 \pm 0.655$ & $+0.477 \pm 0.601$ \\
4v4 & $+0.219 \pm 0.614$ & $+0.438 \pm 0.612$
    & $+0.258 \pm 0.631$ & $+0.367 \pm 0.626$
    & $+0.141 \pm 0.661$ & $+0.375 \pm 0.652$ \\
6v6 & $+0.039 \pm 0.385$ & $+0.102 \pm 0.373$
    & $+0.039 \pm 0.459$ & $+0.039 \pm 0.385$
    & $+0.070 \pm 0.617$ & $-0.047 \pm 0.329$ \\
\hline
\end{tabular}
\end{table*}

At 2v2, all six paired changes from nominal are within \(0.05\) and every
95\% interval of the change includes zero: localization
\((-0.016,+0.039)\), motion \((+0.016,+0.047)\), delay
\((+0.039,-0.016)\) for \((\Delta_A,\Delta_B)\).

At 4v4, specialization stays positive on both poles under every
perturbation, and every interval of the change again includes zero. The
point estimates lean downward on Pole~B: localization
\((-0.016,-0.047)\), motion \((+0.023,-0.117)\), delay \((-0.094,-0.109)\).

At 6v6, localization \((+0.008,+0.008)\) and motion \((+0.008,-0.055)\)
leave both effects positive. Control delay keeps \(\Delta_A\) positive
(\(+0.070\)) but flips \(\Delta_B\) to \(-0.047\), a change of \(-0.141\)
from nominal.

The specialization is therefore stable under localization and motion noise
at every team size, and under control delay at 2v2 and 4v4. Delay is the one
perturbation that removes the Pole-B distinction at 6v6, which suggests that
coordinated defense in larger teams is the component most sensitive to
delayed control.

\paragraph{On subset views.}
All headline numbers above use the sealed \(n=128\) evaluations.
A separate post-hoc descriptive analysis ranks seeds by deployed performance
\(\mathrm{outcome}(A{\to}A)+\mathrm{outcome}(B{\to}B)\) (tie-break: ascending
seed ID), retains the top 50 per method/team size, and for noise selects that
subset from the nominal condition only so perturbations stay paired
(\texttt{true\_top50\_tables.tex},
\texttt{fig\_*\_true\_top50}).
That \(n=50\) view is secondary and must not replace the sealed \(n=128\)
conclusions. An earlier seed-ID prefix of 50 is archived as
\texttt{seedid50\_prefix\_*} and is not the requested top-50 analysis.
